\paragraph{The idea in plain terms.} Look at one digit, say 5, in one band. The finished grid will have exactly one 5
in each of the band's three rows and exactly one in each of its three boxes. So the three 5s pair the rows with the
boxes: each row gets its own box. Some pairings may already be impossible because 5 has been ruled out of that part of
the row. If a particular row--box combination cannot be part of \emph{any} complete pairing, then 5 cannot go in that
triad, even though nothing in the triad itself rules it out. This one idea, checked with a lookup table,
covers locked candidates between rows and boxes, and it narrows a box that has a single place left for the digit down
to a row with a single place left, which the solver then fills. The \emph{band rule} is this check. The \emph{stack rule} is the same check turned sideways: in a stack, the
digit's three placements pair the stack's three columns with its three boxes, and each box of a stack lies in a
different band, so this is a pairing of bands with columns.


% ----------------------------------------------------------------
\subsection{The $3\times 3$ support}

\paragraph{Example.} Take the table of \cref{fig:band}, rows of the band against boxes:
\[
A=\begin{pmatrix} 1 & 1 & 0 \\ 1 & 1 & 0 \\ 1 & 1 & 1 \end{pmatrix}.
\]
Entry $(0,1)$, row 0 with box 1, is in $\mathrm{supp}(A)$: row 1 can then take box 0 and row 2 box 2, a perfect
matching. Entry $(2,0)$, row 2 with box 0, is not: rows 0 and 1 would then both need box 1, since neither has box 2.
So $(2,0)$ is removed, one of the two red entries in \cref{fig:band}(c).


% ----------------------------------------------------------------
\paragraph{The rules in the solver.} The band rule of \cref{alg:band} (line 3) keeps exactly the cells of triads in
$\mathrm{supp}(A^{d,b})$, so by \cref{lem:band} it is sound, and a zero result is a correct contradiction. The stack rule
(\cref{alg:stack}) keeps exactly the column pairs in $\mathrm{supp}(B^{d,s})$ for each stack, so by \cref{lem:stack} it
is sound. Both rules are the matching constraints of Simonis~\cite{simonis2005sudoku}, Section 5.3. The locked-column step of
JCZSolve is the special case of a band with one column in a stack: its row of
$B^{d,s}$ has a single 1, which every perfect matching uses, so the column's entries in the other bands lie on no
perfect matching.



% ----------------------------------------------------------------
The rule costs four table lookups: three 512-entry lookups turn the 27-bit band into the 9-bit matrix $A$, and a
512-entry table maps $A$ to the cells of the triads that lie on a perfect matching. This is the core of the band
solvers of \cref{sec:background}; it covers the consequences of the digit's row and box constraints within the band (locked candidates between rows
and boxes), while placing a digit, and removing the other digits from its cell, is a separate step.

Transposing the argument gives the \textbf{stack rule}. Within a stack, $d$ occurs once in each of its three columns
and once in each of its three boxes, and the boxes of a stack correspond one-to-one to the bands. So $d$'s placement pairs the three bands one-to-one with the stack's three columns, and a (band,
column) pair can be kept only if it lies on a
perfect matching of the $3\times 3$ graph whose edges are the pairs where $d$ still has a candidate. The matrix is
read from each band's 9-bit set of columns in use, so the rule costs three table lookups per digit, and it only has
to run when some band's column set changes. It subsumes the locked-column step of JCZSolve, which removes a column
from the other bands when a box of this band has candidates in that column alone: that is the special case of a
band vertex with degree one. tdoku gets the same inference from its vertical triads.

